\section{Centered bounds, phase prices, and matched discrepancies}
\label{sec:ope-phase-prices}
The raw OPE estimates must hold in every character phase, including phases
of zero measure for the original mixture. An estimate of the mixed-state
expectation alone does not give this. We prove the required transfer by
bounding all even moments of a compressed whole-word operator. The cutoff
argument keeps a matched discrepancy as a single observable throughout;
it never replaces its small price by separate prices for its two terms.

We use the finite-volume planted inequality and the change-of-cell root
estimate established in Section~\ref{ex-centered-estimates}. In the latter,
a centered observable whose root norm is at most $B$ can be moved, with
the corresponding reflection or conjugation, into a coarser cell at a
margin of three of its original cell spans, still with root norm at most
$B$. The proof uses the entire reflected copy orbit and coordinatewise
log convexity. These statements apply to any observable satisfying the
centered interior and root conditions, including a difference of two
observables.

\subsection{A mixed two-scale chessboard estimate}
Let $W_0,W$ be aligned link tilings of a finite torus, with the side of
$W_0$ dividing that of $W$. Write $\mathcal R_W$ for the root norm.
Suppose fine-group observables $F_j$ have centered prices $A_j\geq0$
at a span dividing the side of $W_0$, and their positions $u_j$ lie in
the required three-span margins of $W_0$. Suppose further that
\[
 \operatorname{dist}_\infty(u_j,u_{j'})>\operatorname{side}(W_0)
                  \quad(j\ne j').
\]
For one fixed cell $C$ of $W$, collect the fine-group observables lying
in $C$. Their product is a copy in $C$ of an interior observable $Z_C$,
and
\begin{equation}\label{eq:ope-exterior-root}
 \mathcal R_W(Z_C)\leq\prod_{u_j\in C}A_j.
\end{equation}
Indeed pull each factor back to the home cell of $W$, with the reflection
and conjugation determined by $C$. Now copy this product into every cell
of $W$. The copied positions still lie in the corresponding fine-tiling
margins, since the tilings are aligned and their sides divide. They occupy
distinct cells of $W_0$. Points copied into different cells of $W$ cannot
share a fine cell; within one coarse cell, reflection preserves distances,
and the displayed separation prevents sharing a fine cell. The planted
inequality on $W_0$ therefore bounds the full copied-product expectation
by
\[
 \left(\prod_{u_j\in C}A_j\right)^{\#\mathrm{Cell}(W)}.
\]
Taking the defining root proves \eqref{eq:ope-exterior-root}. Interiority
of the home product follows by multiplying interior factors.

Let a second group $G_i$ have centered prices $B_i\geq0$ at spans dividing
the side of $W$, with positions $v_i$ in the corresponding margins of $W$.
Assume their mutual distances, and their distances to all $u_j$, exceed
$\operatorname{side}(W)$. The $G_i$ then occupy distinct coarse cells,
and none of those cells contains an $F_j$. Apply the planted inequality
on $W$ to the individual $G_i$ and the products $Z_C$. This gives
\begin{equation}\label{eq:ope-mixed-plant}
 \left|\mathbb E\left[\prod_iG_i(\cdot+v_i)
                         \prod_jF_j(\cdot+u_j)\right]\right|
                   \leq\prod_iB_i\prod_jA_j.
\end{equation}
Grouping by cells loses no factors: the disjoint groups of indices
$\{j:u_j\in C\}$ partition the fine-group list. Empty groups contribute
one.

Here are sufficient geometric conditions for placing both groups.
Use depths $b\leq f\leq l\leq n$ and $f\leq q$. The coarse tiling is
at depth $b$, the fine tiling at depth $l$, the centered prices for the
coarse group at depth $q$, and those for the fine group at depth $n$.
If the group sizes are $I_0,J_0$ and
\begin{equation}\label{eq:ope-two-scale-room}
 10I_0a_f<a_b,\qquad 6J_0a_n<a_l,
\end{equation}
there is a common origin for both tilings placing all coarse-group points
at margin $3a_q$ and all fine-group points at margin $3a_n$.
To prove this, first choose the coarse origin coordinate by coordinate.
For each point, origins giving distance less than $5a_f$ from a cell
boundary occupy an interval of length at most $10a_f$ in the circle of
possible origins modulo $a_b$. The first inequality leaves an allowed
origin. Next choose the fine origin modulo $a_l$ avoiding the intervals
of total length at most $6J_0a_n$ forbidden by the fine-group points.
It can be reached from the chosen coarse origin by a shift of coordinate
size at most $a_l\leq a_f$. This shift decreases the coarse margins by
at most $2a_f$ even with the endpoint conventions for link cells, leaving
$3a_f\geq3a_q$. On the finite lattice, if a cell side is $B$ and the desired integer
margin is $r$, a point forbids the stripe residues
$0,\ldots,r-1$ and $B-r,\ldots,B-1$: at most $2r$ residues
of the origin modulo $B$. Summing over the points leaves a permitted
origin whenever their total count is less than $B$. Multiplying
\eqref{eq:ope-two-scale-room} by $a_N^{-1}$ gives precisely these
strict inequalities, first with margin $5a_f/a_N$, then with margin
$3a_n/a_N$. To attain the permitted fine origin, choose a common lattice
shift with $0\leq s_\mu< a_l/a_N$ in each coordinate. The coarse
$5a_f/a_N$ margin loses at most $2a_f/a_N$, leaving the required three
spans. Nested mesh divisibility makes all sides and margins integers.

A Euclidean separation greater than $2a_b$ implies a maximum-coordinate
separation greater than $a_b$ in four dimensions; likewise for $a_l$.
For fixed physical grid points these inequalities persist on the periodic
lattice for all sufficiently large periods: their finitely many coordinate
differences then do not wrap. They persist at all finer cutoffs by the
exact mesh embedding. Thus Euclidean separation at $2a_b$ within and
across the coarse group, and at $2a_l$ within the fine group, suffices
for \eqref{eq:ope-mixed-plant} on the ordered cutoff tail.

\subsection{Even replicas and the norm in the ground algebra}
Let $Z$ be a finite linear combination of products of local fields, at
fixed coordinate-generic grid configurations. Write $A\in\mathcal A$
for the same linear combination of their compressed whole-word
operators. The cutoff version $Z_{N,p}$ may use labels converging in their
fixed finite-dimensional label spaces along a nontrivial refinement of
the ordered cutoff limit.

Suppose each cutoff $Z_{N,p}$ factors as the product of two groups as in
\eqref{eq:ope-mixed-plant}, with the same centered price bounds and
geometry. The statement includes a group consisting of a single
matched discrepancy. Set $C=\prod_iB_i\prod_jA_j$. For each $m\geq1$,
translate $2m$ copies of $Z_{N,p}$ by the even coarse-cell displacements
\[
 d_j(t)=2tj\,a_b(1,1,1,1),\qquad t\in\mathbb N,\quad 0\leq j<2m,
\]
conjugating the copies with even $j$. These translations preserve all
tilings' margins, because they are multiples of the coarse cell side,
and this side is a multiple of the fine side. For fixed $m$, copies at
distinct indices become separated in every coordinate as $t\to\infty$.
For each sufficiently large fixed $t$, take the period large enough to
avoid wrapping, and then the cutoff fine enough to use all the centered
prices. Applying \eqref{eq:ope-mixed-plant} to the replicated lists gives
\begin{equation}\label{eq:ope-replica-price}
 \left|\mathbb E_{N,p}\prod_{j=0}^{2m-1}
       \begin{cases}\overline{Z_{N,p}(\cdot+d_j(t))},&j\text{ even},\\
                       Z_{N,p}(\cdot+d_j(t)),&j\text{ odd}
       \end{cases}\right|\leq C^{2m}.
\end{equation}
Only the two original arities occur in \eqref{eq:ope-two-scale-room}.
Once an origin works for one copy, it works for all its even-cell
translates. The period and the translation threshold may depend on $m$;
the bound $C$ does not.

Expand the finite polynomial in \eqref{eq:ope-replica-price} and first
pass to the ordered cutoff limit at fixed $t$. For sufficiently large
$t$ each expanded word is coordinate-generic, so the point-limit theorem
applies. Converging labels cause no extra difficulty: in a fixed basis
a multilinear correlation is a finite sum of basis correlations times
products of label coordinates. The basis correlations converge and the
coordinates converge, proving convergence of each term. This also works
on any finer nontrivial cutoff filter. It does not assume that the moving
labels have separate limits in the two cutoff indices.

Next let $t\to\infty$. The separated-cluster limit of
Lemma~\ref{ex-cluster-train}, together with conjugation becoming adjoint
under ground compression, identifies the limit as
\[
 \langle\Omega,(A^*A)^m\Omega\rangle.
\]
For a linear combination this follows by expanding the ordered product;
its alternating factors are exactly $A^*,A,\ldots,A^*,A$. Thus
\[
 \int_\Xi|\chi(A)|^{2m}\,d\mu(\chi)\leq C^{2m}
                  \quad(m\geq1).
\]
Full support of $\mu$ now gives $|\chi(A)|\leq C$ for every character.
If one value exceeded $C$, continuity would give an open set of positive
measure on which it exceeded some $C'>C$. Then
$\mu(O)(C')^{2m}\leq C^{2m}$ for every $m$, a contradiction. The same
argument covers $C=0$. The isometric Gelfand transform consequently gives
\begin{equation}\label{eq:ope-ground-price}
 \|A\|\leq C.
\end{equation}
Both every character evaluation and the mixed-state functional have norm
one, so the same bound holds for their evaluations of $A$.

\subsection{Removing coordinate ties}
To extend \eqref{eq:ope-ground-price} to arbitrary separated
configurations, first approximate all their points by the generic grids
of \eqref{eq:ope-grid-round}. Strict Euclidean separations persist.
For each fixed character, its correlation is continuous on separated
configurations, so the generic bound passes to the limit. Taking the
supremum over characters then proves the ground-algebra norm bound at
the original configuration. This uses scalar continuity in each phase,
not norm continuity of the ground-algebra-valued correlation.

For a matched discrepancy whose cluster is already fixed and generic,
only the spectator positions are approximated. They must remain
coordinate-generic relative to both that cluster and the origin. Here is
the required extension of the grid construction. Let $W$ be the finite
set consisting of the cluster points and the origin, $J_0$ the number of
spectators, and $K=(J_0+1)(|W|+1)$. In each coordinate choose for spectator
$j$ a residue from the block
\[
 \{j(|W|+1),\ldots,j(|W|+1)+|W|\}
\]
which avoids the residues $\lfloor w_\mu/h\rfloor\bmod K$ for $w\in W$.
There are $|W|+1$ candidates and at most $|W|$ forbidden residues. The
blocks are disjoint, so the chosen residues distinguish all spectators.
Adding the chosen residue to
$K\lfloor y_{j,\mu}/(hK)\rfloor$ gives a grid point within $hK$ in
each coordinate, and it cannot equal any $w_\mu$. Thus these spectator
grids converge and retain all required coordinate inequalities. If the
new grid is finer than the fixed cluster grid, reexpress the cluster
points on it by the integer mesh ratio; their physical positions and
cutoff observables are unchanged. Apply the preceding replica argument
on that common grid, then take the spectator approximation limit in
each phase. This proves the complete-discrepancy bound at arbitrary
separated spectator configurations.

\subsection{Engineering rates from centered prices}
We now use the centered estimate in its full quantifier order: for a
fixed degree $d$, bounded set of labels and every $\epsilon>0$, one
constant controls the centered root price at all depths after the
reference depth, with an eventual cutoff tail uniform in period and
those labels. In particular a fixed field of degree $d_i$ has a price
\[
 B_i(k)=(1+C_i)a_k^{-d_i-\epsilon}+1
\]
for a suitable $C_i\geq0$. If $s_{k_n}/t_n\to u\geq0$, this positive
price has logarithmic rate $(d_i+\epsilon)u$. For a finite group the
product therefore has rate $(\sum_i d_i+\epsilon\#I)u$.

More generally, suppose two groups have mutual separations at upper
scales $0\leq a\leq b$, and cross separation at scale $a$. This means that
for every positive slack the corresponding distances eventually exceed
$e^{-(a+\delta)t_n}$ or $e^{-(b+\delta)t_n}$. Choose mesh depths of
logarithmic rates
\[
 a+\delta/3,\quad a+\delta/2,\quad a+\delta,
 \quad b+3\delta/4,\quad b+\delta
\]
for the coarse placement, buffer, coarse-group price, fine placement,
and fine-group price respectively. The inequalities of the rates imply
all depth orderings and \eqref{eq:ope-two-scale-room} eventually. The
separation assumptions imply the required strict gaps at the placement
meshes. The bounds just proved yield ground-algebra upper rate
\[
 (w_1+\epsilon I_0)(a+\delta)
                  +(w_2+\epsilon J_0)(b+\delta).
\]
Let $\epsilon,\delta\downarrow0$ by allocating half of any desired
rate slack to these choices and half to the eventual estimate. The
result is the upper rate $w_1a+w_2b$, uniformly along every sequence of
phases, because the ground-algebra norm controls them all.

For the OPE, take the fixed-clearance spectators as the first group
($a=0$) and the scale-$x$ cluster as the second ($b=x$).
Their cross clearance is fixed for all sufficiently large indices.
This proves the rough rate $xw(p)$ used in Section~\ref{sec:ope-grid}.
Taking every group at scale zero proves rate zero for the fixed
single-field spectator insertion in Input~\ref{input:ope-raw-estimates}.

Finally suppose a cutoff matched discrepancy is centered and priced by
$A_n$ at a depth of logarithmic rate $b>0$, and its coefficient label
converges along a refinement of the same ordered cutoff limit. Assume
$A_n$ has upper rate $\gamma$. Choose, for any small $0<\delta<b$,
placement, buffer, fine-placement and spectator-price depths of rates
$\delta/4,\delta/2,3\delta/4,\delta$. The centered discrepancy price
at rate $b$ fits this geometry. Equation~\eqref{eq:ope-ground-price}
and the spectator approximation give upper rate at most
$\gamma+(w_2+\epsilon J_0)\delta$ for the entire discrepancy, in every
phase and at arbitrary fixed-clearance spectator positions. Sending
$\delta\downarrow0$ proves rate $\gamma$. Thus the finite-cutoff match
supplies exactly the complete-discrepancy part of
Input~\ref{input:ope-match}, once its price and convergent label are
constructed. Neither this transfer nor the even-moment argument creates
that finite-cutoff match.
